\section{讨论与后续研究方向}\label{ch:open-questions}

\subsection{从最小化训练风险到控制模型复杂度}

前述方法主要研究给定模型族内的参数求解。对于模型选择和泛化能力，仍需要考察正则化的作用。一个自然的扩展是
\begin{equation}\label{eq:regularized-risk}
\min_{\theta}\left\{\frac1m\sum_{i=1}^{m}L\bigl(y_i,f_\theta(x_i)\bigr)
+\lambda\Omega(\theta)\right\},\qquad\lambda\ge0.
\end{equation}
$\Omega(\theta)$ 用于表达偏好的参数结构或复杂度，$\lambda$ 调节经验拟合与附加约束之间的权衡。岭回归、Lasso 及其相关优化条件是这一方向的后续重点，现阶段不将尚未讨论的收敛性或泛化效果当作既定结论。

\subsection{从解析方法到一般数值优化}

最小二乘目标的 Hessian 为 $X\trans X/m$，其凸性和条件数会直接影响数值求解。今后应进一步比较正规方程、QR/SVD 分解、梯度下降、线搜索和牛顿类方法：哪些方法依赖凸性，哪些方法能处理不可微目标，何时需要考虑约束，以及计算规模如何影响算法选择。对于具备二次规划结构的问题，还可以适当考察专用优化求解器的作用，但不应以求解器替代数学证明。

\subsection{从点估计到统计推断}

最小二乘系数作为点估计后，还可进一步研究其偏差、方差与区间估计。相关结论需要区分模型形式、误差条件、设计矩阵秩条件以及是否将输入视为随机变量。尤其应避免把“估计系数显著”与“预测效果好”直接等同。

\subsection{仍需回答的问题}

\begin{enumerate}
\item 对同一个预测任务，不同损失函数分别隐含怎样的统计假设和决策偏好？
\item 对于凸但不强凸的目标，梯度下降在何种条件下有函数值收敛率？
\item 当样本数接近或小于特征数时，最小二乘估计的不唯一性如何影响泛化？
\item 当训练与测试数据并非来自同一分布时，现有交叉验证流程需要怎样调整？
\end{enumerate}

这些问题将在后续专题中分别展开，而不是在本节提前给出没有证明条件的统一回答。
